% =====================================================================
\section{Threats to Validity}\label{sec:threats}
\emph{Construct validity.} Properties P1--P7 and L1 operationalize Theorems~\ref{thm:safety}, \ref{thm:cost} and~\ref{thm:termination}. The RQ1 oracle recomputes observations from the specification, but it shares the case study's definitions of read set and validator with the controller, so a shared misunderstanding of those definitions would go undetected. Request counts measure avoided generation calls. In RQ5 we also report tokens and wall-clock time, but these depend on hardware, quantization and server settings and should be read comparatively. Our equality categories for early cutoff (raw, value, AST, behavior) are proxies. Behavioral equality on the plus inputs is an under-approximation of semantic equivalence, and AST equality is an over-approximation of what an exact key may safely treat as equal unless the normal form is part of the contract.

\emph{Internal validity.} Each RQ1 ablation reuses the reference code and overrides only the removed check. During development we found and fixed one harness problem, a base-class candidate constructor that silently bypassed overrides. Scenario expectations were written by the author before execution and are conformance tests, not independent evidence. In RQ5, the prompts, clauses and validators were written by the author. Two clarifications were made after a smoke test with one model: archived documents remain readable, and tests may use only public methods. They removed ambiguities that made correct programs fail, and they were fixed before any reported run. One ambiguity remained undetected until the analysis: the design contract does not name the operations, while the validator requires the name \texttt{put} used in the clauses, and most first design attempts used \texttt{update} (Section~\ref{sec:rq5}). This depresses first-attempt acceptance of designs and inflates the number of retries, but it cannot affect the safety outcomes, which hold for any proposal. We report the runs as executed. The stale-reply trials flip R03 half a second after dispatch. Every reply therefore arrives after the change, which isolates the question ``how often is a late reply wrong?'' from the question of how often replies are late in a real deployment.

\emph{External validity.} The scripted case study is small, and the RQ2 workspaces are synthetic. The coupled workspace adds shared components with controlled fan-out, but real software has richer coupling. We planned to include \texttt{deepseek-coder-v2:16b}, but on our Ollama build it needed minutes per short completion even when fully GPU-resident, so we replaced it with \texttt{gemma2:9b} before collecting any reported data. RQ5 uses one realistic but small workflow with seven open models of 7--30 billion parameters and six workspaces per model, so per-model rates are anecdotal and we report pooled proportions with Wilson 95\% intervals. Frontier models may fail less often; the coordination guarantees do not depend on this, but the frequencies do. RQ6 covers two function-level benchmarks. The validator-strength gap we measure is specific to their base and plus tests, and repository-level tasks such as SWE-bench~\cite{swebench} would need structural proposals (Section~\ref{sec:limitations}). Bounded model checking covers instances of up to \TLCExtMaxJobs{} jobs.

\emph{Conclusion validity.} RQ1--RQ3 runs are deterministic functions of their seeds and are not samples from a deployment population, so we report counts and proportions, not significance tests. In RQ5 and RQ6, model outputs are seeded per request but GPU inference is not bit-reproducible across hardware. We therefore log every response and derive every number from the logs. With five to seven models and a few hundred workflow runs, we report descriptive statistics and avoid claims about differences between individual models.
